\section{Bài 6}

\subsection{Đề bài}

Vẽ mạch cụ thể và kiểm tra tính đúng đắn của thuật toán Deutsch-Jozsa cho các trường hợp hàm $f_1$ và $f_2$ ở bảng sau

\begin{table}[H]
    \centering
    \begin{tabular}{cc|c|c}
        \hline
        $x_1$ & $x_0$ & $f_1$ & $f_2$ \\
        \hline
        0 & 0 & 1 & 0 \\
        0 & 1 & 1 & 0 \\
        1 & 0 & 1 & 1 \\
        1 & 1 & 1 & 1 \\
        \hline
    \end{tabular}
\end{table}

\subsection{Lời giải}

Ta xây dựng hàm lượng tử $U_{f_1}$ và $U_{f_2}$ tương ứng với hai hàm $f_1$ và $f_2$ như sau:
\begin{itemize}
    \item Hàm $f_1$ là hàm hằng, ta có thể xây dựng $U_{f_1}$ bằng cách áp dụng cổng X lên qubit thứ 2 (qubit mục tiêu).
    \item Hàm $f_2$ là hàm cân bằng, ta có thể xây dựng $U_{f_2}$ bằng cách áp dụng cổng CNOT với qubit thứ 1 làm điều khiển và qubit thứ 2 làm mục tiêu.
\end{itemize}

Mạch Deutsch-Jozsa cho $f_1$: 
\begin{center}
    \begin{quantikz}
    \lstick{$\ket{0} (x_1)$} & \gate{H} & \qw       & \gate{H} & \meter{} & \rstick{$\ket{0}$} \qw \\
    \lstick{$\ket{0} (x_0)$} & \gate{H} & \qw       & \gate{H} & \meter{} & \rstick{$\ket{0}$} \qw \\
    \lstick{$\ket{1} (y)$}   & \gate{H} & \gate{X}  & \qw      & \qw      & \qw
    \end{quantikz}
\end{center}

Mạch Deutsch-Jozsa cho $f_2$:
\begin{center}
    \begin{quantikz}
        \lstick{$\ket{0} (x_1)$} & \gate{H} & \ctrl{2} & \gate{H} & \meter{} & \rstick{$\ket{1}$} \qw \\
        \lstick{$\ket{0} (x_0)$} & \gate{H} & \qw      & \gate{H} & \meter{} & \rstick{$\ket{0}$} \qw \\
        \lstick{$\ket{1} (y)$}   & \gate{H} & \targ{}  & \qw      & \qw      & \qw
    \end{quantikz}
\end{center}

Sau khi thực hiện thuật toán Deutsch-Jozsa, ta đo hai qubit đầu tiên. Kết quả đo được sẽ là:
\begin{itemize}
    \item Đối với hàm $f_1$ (hàm hằng), ta sẽ luôn đo được kết quả $\ket{00}$.
    \item Đối với hàm $f_2$ (hàm cân bằng), ta sẽ đo được kết quả khác $\ket{10}$.
\end{itemize}

